
\documentclass{article}
\usepackage{colortbl}
\usepackage{makecell}
\usepackage{multirow}
\usepackage{supertabular}

\begin{document}

\newcounter{utterance}

\centering \large Interaction Transcript for game `clean\_up', experiment `0\_easy\_3obj\_es', episode 4 with Apertus{-}v1.5{-}8B.
\vspace{24pt}

{ \footnotesize  \setcounter{utterance}{1}
\setlength{\tabcolsep}{0pt}
\begin{supertabular}{c@{$\;$}|p{.15\linewidth}@{}p{.15\linewidth}p{.15\linewidth}p{.15\linewidth}p{.15\linewidth}p{.15\linewidth}}
    \# & $\;$A & \multicolumn{4}{c}{Game Master} & $\;\:$B\\
    \hline

    \theutterance \stepcounter{utterance}  
    & & \multicolumn{4}{p{0.6\linewidth}}{
        \cellcolor[rgb]{0.9,0.9,0.9}{
            \makecell[{{p{\linewidth}}}]{
                \texttt{\tiny{[P1$\langle$GM]}}
                \texttt{Soy tu maestro de juego, y estás jugando un juego colaborativo con la siguiente cuadrícula como tablero de juego:} \\
\\ 
\texttt{```} \\
\texttt{ 1234567} \\
\texttt{╔═╤═════╗} \\
\texttt{║◌│◌◌◌◌L║ 1} \\
\texttt{║◌│◌◌◌◌◌║ 2} \\
\texttt{║◌│◌◌◌◌◌║ 3} \\
\texttt{║◌│◌◌◌◌◌║ 4} \\
\texttt{╟─┼───┐C║ 5} \\
\texttt{║◌│◌◌◌├─╢ 6} \\
\texttt{║P│◌◌◌│◌║ 7} \\
\texttt{╚═╧═══╧═╝} \\
\\ 
\texttt{```} \\
\\ 
\texttt{* El borde superior muestra coordenadas x que aumentan hacia la derecha, y el borde derecho coordenadas y que aumentan hacia abajo.} \\
\texttt{* Los siguientes objetos están colocados aleatoriamente en tu cuadrícula: 'C', 'L', 'P'.} \\
\\ 
\\ 
\texttt{El otro jugador ve una variación del tablero de juego, donde los objetos están colocados en diferentes ubicaciones aleatorias. No puedes ver el tablero del otro jugador, y él no puede ver el tuyo.} \\
\\ 
\\ 
\texttt{**Objetivo:**} \\
\\ 
\texttt{Ambos jugadores necesitan mover los objetos en sus respectivos fondos de manera que los objetos idénticos terminen en las mismas coordenadas. Tienes que comunicarte con el otro jugador para acordar una configuración objetivo común.} \\
\\ 
\\ 
\texttt{**Reglas:**} \\
\\ 
\texttt{* En cada turno, puedes enviar exactamente uno de los dos comandos siguientes:} \\
\texttt{ 1. `DECIR: <MENSAJE>`: para enviar un mensaje (todo hasta el siguiente salto de línea) al otro jugador. Se lo reenviaré a tu compañero.} \\
\texttt{2. `MOVERR: <OBJETO>, (<X>, <Y>)`: para mover un objeto a una nueva posición, donde `<X>` es la columna y `<Y>` es la fila. Te informaré si tu movimiento fue exitoso o no.} \\
\texttt{* Si no sigues el formato, o envías varios comandos a la vez, tendré que penalizarte.} \\
\texttt{* Si ambos jugadores acumulan más de 8 penalizaciones, ambos pierden el juego.} \\
\texttt{* ¡Es vital que te comuniques con el otro jugador respecto a tu estado objetivo! ¡La *única* forma de transmitir tu estrategia al otro jugador es usando el comando `DECIR: <MENSAJE>`!} \\
\\ 
\\ 
\texttt{**Mover objetos**} \\
\\ 
\texttt{* Solo puedes mover objetos a celdas dentro de los límites de la cuadrícula. La celda de destino debe estar vacía, es decir, solo debe contener el símbolo '◌'.} \\
\texttt{* Si intentas mover un objeto a un lugar que no está vacío, o intentas moverlo fuera de la cuadrícula, tendré que penalizarte. Tendrás otro intento.} \\
\texttt{* ¡Antes de hacer un movimiento, verifica dos veces que el lugar de destino esté vacío, y no contenga ninguna letra, marco o línea!} \\
\\ 
\\ 
\texttt{**Fin del juego**} \\
\\ 
\texttt{Si crees que alcanzaste el objetivo de alinear todos los objetos, puedes pedirle al otro jugador que termine el juego enviando `DECIR: ¿terminado?`. Si el otro jugador te pide que termines el juego, y respondes con `DECIR: ¡terminado!`, el juego terminará.} \\
\\ 
\texttt{Ambos jugadores ganan si el juego termina dentro de 12 rondas, donde una ronda se define como dos jugadores enviando cada uno un comando válido.} \\
\\ 
\texttt{**Puntuación:**} \\
\\ 
\texttt{Cuanto más cerca estén los objetos idénticos en ambos tableros de juego, más puntos obtendrás. Las penalizaciones reducen tus puntos. ¿Puedes batir el récord?} \\
\\ 
\\ 
\texttt{¡Por favor, envía un mensaje al otro jugador para iniciar el juego!} \\
            }
        }
    }
    & \\ \\

    \theutterance \stepcounter{utterance}  
    & \multicolumn{4}{p{0.6\linewidth}}{
        \cellcolor[rgb]{0.9,0.9,0.9}{
            \makecell[{{p{\linewidth}}}]{
                \texttt{\tiny{[P1$\rangle$GM]}}
                \texttt{DECIR: Hola, vamos a coordinar nuestros objetivos. ¿Podrías confirmar la ubicación de tu 'C', 'L' y 'P'? Me gustaría establecer una configuración común.} \\
            }
        }
    }
    & & \\ \\

    \theutterance \stepcounter{utterance}  
    & & & \multicolumn{2}{p{0.3\linewidth}}{
        \cellcolor[rgb]{0.9,0.9,0.9}{
            \makecell[{{p{\linewidth}}}]{
                \texttt{\tiny{[GM$|$GM]}}
                \texttt{DECIR: Hola, vamos a coordinar nuestros objetivos. ¿Podrías confirmar la ubicación de tu 'C', 'L' y 'P'? Me gustaría establecer una configuración común.} \\
            }
        }
    }
    & & \\ \\

    \theutterance \stepcounter{utterance}  
    & & \multicolumn{4}{p{0.6\linewidth}}{
        \cellcolor[rgb]{0.9,0.9,0.9}{
            \makecell[{{p{\linewidth}}}]{
                \texttt{\tiny{[GM$\rangle$P2]}}
                \texttt{Soy tu maestro de juego, y estás jugando un juego colaborativo con la siguiente cuadrícula como tablero de juego:} \\
\\ 
\texttt{```} \\
\texttt{ 1234567} \\
\texttt{╔═╤═════╗} \\
\texttt{║◌│◌◌◌◌◌║ 1} \\
\texttt{║◌│◌◌PL◌║ 2} \\
\texttt{║◌│◌◌◌◌◌║ 3} \\
\texttt{║◌│◌◌◌◌◌║ 4} \\
\texttt{╟─┼───┐◌║ 5} \\
\texttt{║◌│◌◌◌├─╢ 6} \\
\texttt{║◌│◌C◌│◌║ 7} \\
\texttt{╚═╧═══╧═╝} \\
\\ 
\texttt{```} \\
\\ 
\texttt{* El borde superior muestra coordenadas x que aumentan hacia la derecha, y el borde derecho coordenadas y que aumentan hacia abajo.} \\
\texttt{* Los siguientes objetos están colocados aleatoriamente en tu cuadrícula: 'C', 'L', 'P'.} \\
\\ 
\\ 
\texttt{El otro jugador ve una variación del tablero de juego, donde los objetos están colocados en diferentes ubicaciones aleatorias. No puedes ver el tablero del otro jugador, y él no puede ver el tuyo.} \\
\\ 
\\ 
\texttt{**Objetivo:**} \\
\\ 
\texttt{Ambos jugadores necesitan mover los objetos en sus respectivos fondos de manera que los objetos idénticos terminen en las mismas coordenadas. Tienes que comunicarte con el otro jugador para acordar una configuración objetivo común.} \\
\\ 
\\ 
\texttt{**Reglas:**} \\
\\ 
\texttt{* En cada turno, puedes enviar exactamente uno de los dos comandos siguientes:} \\
\texttt{ 1. `DECIR: <MENSAJE>`: para enviar un mensaje (todo hasta el siguiente salto de línea) al otro jugador. Se lo reenviaré a tu compañero.} \\
\texttt{2. `MOVERR: <OBJETO>, (<X>, <Y>)`: para mover un objeto a una nueva posición, donde `<X>` es la columna y `<Y>` es la fila. Te informaré si tu movimiento fue exitoso o no.} \\
\texttt{* Si no sigues el formato, o envías varios comandos a la vez, tendré que penalizarte.} \\
\texttt{* Si ambos jugadores acumulan más de 8 penalizaciones, ambos pierden el juego.} \\
\texttt{* ¡Es vital que te comuniques con el otro jugador respecto a tu estado objetivo! ¡La *única* forma de transmitir tu estrategia al otro jugador es usando el comando `DECIR: <MENSAJE>`!} \\
\\ 
\\ 
\texttt{**Mover objetos**} \\
\\ 
\texttt{* Solo puedes mover objetos a celdas dentro de los límites de la cuadrícula. La celda de destino debe estar vacía, es decir, solo debe contener el símbolo '◌'.} \\
\texttt{* Si intentas mover un objeto a un lugar que no está vacío, o intentas moverlo fuera de la cuadrícula, tendré que penalizarte. Tendrás otro intento.} \\
\texttt{* ¡Antes de hacer un movimiento, verifica dos veces que el lugar de destino esté vacío, y no contenga ninguna letra, marco o línea!} \\
\\ 
\\ 
\texttt{**Fin del juego**} \\
\\ 
\texttt{Si crees que alcanzaste el objetivo de alinear todos los objetos, puedes pedirle al otro jugador que termine el juego enviando `DECIR: ¿terminado?`. Si el otro jugador te pide que termines el juego, y respondes con `DECIR: ¡terminado!`, el juego terminará.} \\
\\ 
\texttt{Ambos jugadores ganan si el juego termina dentro de 12 rondas, donde una ronda se define como dos jugadores enviando cada uno un comando válido.} \\
\\ 
\texttt{**Puntuación:**} \\
\\ 
\texttt{Cuanto más cerca estén los objetos idénticos en ambos tableros de juego, más puntos obtendrás. Las penalizaciones reducen tus puntos. ¿Puedes batir el récord?} \\
\\ 
\texttt{El otro jugador inició el juego enviando este mensaje:} \\
\\ 
\texttt{"Hola, vamos a coordinar nuestros objetivos. ¿Podrías confirmar la ubicación de tu 'C', 'L' y 'P'? Me gustaría establecer una configuración común."} \\
\\ 
\texttt{¿Cuál es tu primer comando?} \\
            }
        }
    }
    & \\ \\

    \theutterance \stepcounter{utterance}  
    & & & \multicolumn{4}{p{0.6\linewidth}}{
        \cellcolor[rgb]{0.9,0.9,0.9}{
            \makecell[{{p{\linewidth}}}]{
                \texttt{\tiny{[GM$\langle$P2]}}
                \texttt{DECIR: Estoy analizando mi tablero. Mi 'C' está en la celda (7,5), mi 'L' está en la celda (2,2) y mi 'P' está en la celda (2,3). Mi objetivo es coordinar estas ubicaciones con las tuyas para lograr que todos los objetos idénticos estén alineados en posiciones comunes. ¿Podrías compartir la ubicación de tu 'C', 'L' y 'P'?} \\
            }
        }
    }
     \\ \\

    \theutterance \stepcounter{utterance}  
    & & & \multicolumn{2}{p{0.3\linewidth}}{
        \cellcolor[rgb]{0.9,0.9,0.9}{
            \makecell[{{p{\linewidth}}}]{
                \texttt{\tiny{[GM$|$GM]}}
                \texttt{DECIR: Estoy analizando mi tablero. Mi 'C' está en la celda (7,5), mi 'L' está en la celda (2,2) y mi 'P' está en la celda (2,3). Mi objetivo es coordinar estas ubicaciones con las tuyas para lograr que todos los objetos idénticos estén alineados en posiciones comunes. ¿Podrías compartir la ubicación de tu 'C', 'L' y 'P'?} \\
            }
        }
    }
    & & \\ \\

    \theutterance \stepcounter{utterance}  
    & & \multicolumn{4}{p{0.6\linewidth}}{
        \cellcolor[rgb]{0.9,0.9,0.9}{
            \makecell[{{p{\linewidth}}}]{
                \texttt{\tiny{[P1$\langle$GM]}}
                \texttt{Tu mensaje ha sido transmitido al otro jugador.} \\
\\ 
\\ 
\texttt{Actualmente estás jugando la ronda 1 de un máximo de 12.} \\
\texttt{Han acumulado colectivamente 0 de 8 penalizaciones.} \\
\texttt{El otro jugador envió este mensaje:} \\
\\ 
\\ 
\texttt{"Estoy analizando mi tablero. Mi 'C' está en la celda (7,5), mi 'L' está en la celda (2,2) y mi 'P' está en la celda (2,3). Mi objetivo es coordinar estas ubicaciones con las tuyas para lograr que todos los objetos idénticos estén alineados en posiciones comunes. ¿Podrías compartir la ubicación de tu 'C', 'L' y 'P'?"} \\
\\ 
\\ 
\texttt{¿Cuál es tu próximo comando?} \\
            }
        }
    }
    & \\ \\

    \theutterance \stepcounter{utterance}  
    & \multicolumn{4}{p{0.6\linewidth}}{
        \cellcolor[rgb]{0.9,0.9,0.9}{
            \makecell[{{p{\linewidth}}}]{
                \texttt{\tiny{[P1$\rangle$GM]}}
                \texttt{DECIR: Mi 'C' está en (5,5), mi 'L' está en (1,2) y mi 'P' está en (2,1). Vamos a intentar alinear nuestras posiciones para que coincidan. ¿Qué te parece si movemos nuestras 'L' hacia (2,2) y nuestra 'P' hacia (2,3) para emparejarlas con las tuyas? ¿Estás de acuerdo con este plan?} \\
            }
        }
    }
    & & \\ \\

    \theutterance \stepcounter{utterance}  
    & & & \multicolumn{2}{p{0.3\linewidth}}{
        \cellcolor[rgb]{0.9,0.9,0.9}{
            \makecell[{{p{\linewidth}}}]{
                \texttt{\tiny{[GM$|$GM]}}
                \texttt{DECIR: Mi 'C' está en (5,5), mi 'L' está en (1,2) y mi 'P' está en (2,1). Vamos a intentar alinear nuestras posiciones para que coincidan. ¿Qué te parece si movemos nuestras 'L' hacia (2,2) y nuestra 'P' hacia (2,3) para emparejarlas con las tuyas? ¿Estás de acuerdo con este plan?} \\
            }
        }
    }
    & & \\ \\

    \theutterance \stepcounter{utterance}  
    & & \multicolumn{4}{p{0.6\linewidth}}{
        \cellcolor[rgb]{0.9,0.9,0.9}{
            \makecell[{{p{\linewidth}}}]{
                \texttt{\tiny{[GM$\rangle$P2]}}
                \texttt{Tu mensaje ha sido transmitido al otro jugador.} \\
\\ 
\\ 
\texttt{Actualmente estás jugando la ronda 2 de un máximo de 12.} \\
\texttt{Han acumulado colectivamente 0 de 8 penalizaciones.} \\
\texttt{El otro jugador envió este mensaje:} \\
\\ 
\\ 
\texttt{"Mi 'C' está en (5,5), mi 'L' está en (1,2) y mi 'P' está en (2,1). Vamos a intentar alinear nuestras posiciones para que coincidan. ¿Qué te parece si movemos nuestras 'L' hacia (2,2) y nuestra 'P' hacia (2,3) para emparejarlas con las tuyas? ¿Estás de acuerdo con este plan?"} \\
\\ 
\\ 
\texttt{¿Cuál es tu próximo comando?} \\
            }
        }
    }
    & \\ \\

    \theutterance \stepcounter{utterance}  
    & & & \multicolumn{4}{p{0.6\linewidth}}{
        \cellcolor[rgb]{0.9,0.9,0.9}{
            \makecell[{{p{\linewidth}}}]{
                \texttt{\tiny{[GM$\langle$P2]}}
                \texttt{DECIR: Estoy de acuerdo con el plan propuesto. Para alinear nuestras 'L' y 'P', moveré mi 'L' desde (2,2) a (2,2) (mis coordenadas actuales) y mi 'P' desde (2,3) a (2,3) (mis coordenadas actuales). Espero que esto ayude a emparejar nuestras posiciones. ¿Puedes confirmar que tu 'L' en (2,2) y 'P' en (2,3) no hayan sido ocupadas por otros objetos al momento de moverlas?} \\
            }
        }
    }
     \\ \\

    \theutterance \stepcounter{utterance}  
    & & & \multicolumn{2}{p{0.3\linewidth}}{
        \cellcolor[rgb]{0.9,0.9,0.9}{
            \makecell[{{p{\linewidth}}}]{
                \texttt{\tiny{[GM$|$GM]}}
                \texttt{DECIR: Estoy de acuerdo con el plan propuesto. Para alinear nuestras 'L' y 'P', moveré mi 'L' desde (2,2) a (2,2) (mis coordenadas actuales) y mi 'P' desde (2,3) a (2,3) (mis coordenadas actuales). Espero que esto ayude a emparejar nuestras posiciones. ¿Puedes confirmar que tu 'L' en (2,2) y 'P' en (2,3) no hayan sido ocupadas por otros objetos al momento de moverlas?} \\
            }
        }
    }
    & & \\ \\

    \theutterance \stepcounter{utterance}  
    & & \multicolumn{4}{p{0.6\linewidth}}{
        \cellcolor[rgb]{0.9,0.9,0.9}{
            \makecell[{{p{\linewidth}}}]{
                \texttt{\tiny{[P1$\langle$GM]}}
                \texttt{Tu mensaje ha sido transmitido al otro jugador.} \\
\\ 
\\ 
\texttt{Actualmente estás jugando la ronda 2 de un máximo de 12.} \\
\texttt{Han acumulado colectivamente 0 de 8 penalizaciones.} \\
\texttt{El otro jugador envió este mensaje:} \\
\\ 
\\ 
\texttt{"Estoy de acuerdo con el plan propuesto. Para alinear nuestras 'L' y 'P', moveré mi 'L' desde (2,2) a (2,2) (mis coordenadas actuales) y mi 'P' desde (2,3) a (2,3) (mis coordenadas actuales). Espero que esto ayude a emparejar nuestras posiciones. ¿Puedes confirmar que tu 'L' en (2,2) y 'P' en (2,3) no hayan sido ocupadas por otros objetos al momento de moverlas?"} \\
\\ 
\\ 
\texttt{¿Cuál es tu próximo comando?} \\
            }
        }
    }
    & \\ \\

    \theutterance \stepcounter{utterance}  
    & \multicolumn{4}{p{0.6\linewidth}}{
        \cellcolor[rgb]{0.9,0.9,0.9}{
            \makecell[{{p{\linewidth}}}]{
                \texttt{\tiny{[P1$\rangle$GM]}}
                \texttt{DECIR: Confirmo que mi 'L' está en (2,2) y mi 'P' está en (2,3), y en ambos casos las celdas están vacías, por lo que no hay objetos que impidan el movimiento. ¿Estás listo para realizar los movimientos correspondientes en tu tablero?} \\
            }
        }
    }
    & & \\ \\

    \theutterance \stepcounter{utterance}  
    & & & \multicolumn{2}{p{0.3\linewidth}}{
        \cellcolor[rgb]{0.9,0.9,0.9}{
            \makecell[{{p{\linewidth}}}]{
                \texttt{\tiny{[GM$|$GM]}}
                \texttt{DECIR: Confirmo que mi 'L' está en (2,2) y mi 'P' está en (2,3), y en ambos casos las celdas están vacías, por lo que no hay objetos que impidan el movimiento. ¿Estás listo para realizar los movimientos correspondientes en tu tablero?} \\
            }
        }
    }
    & & \\ \\

    \theutterance \stepcounter{utterance}  
    & & \multicolumn{4}{p{0.6\linewidth}}{
        \cellcolor[rgb]{0.9,0.9,0.9}{
            \makecell[{{p{\linewidth}}}]{
                \texttt{\tiny{[GM$\rangle$P2]}}
                \texttt{Tu mensaje ha sido transmitido al otro jugador.} \\
\\ 
\\ 
\texttt{Actualmente estás jugando la ronda 3 de un máximo de 12.} \\
\texttt{Han acumulado colectivamente 0 de 8 penalizaciones.} \\
\texttt{El otro jugador envió este mensaje:} \\
\\ 
\\ 
\texttt{"Confirmo que mi 'L' está en (2,2) y mi 'P' está en (2,3), y en ambos casos las celdas están vacías, por lo que no hay objetos que impidan el movimiento. ¿Estás listo para realizar los movimientos correspondientes en tu tablero?"} \\
\\ 
\\ 
\texttt{¿Cuál es tu próximo comando?} \\
            }
        }
    }
    & \\ \\

    \theutterance \stepcounter{utterance}  
    & & & \multicolumn{4}{p{0.6\linewidth}}{
        \cellcolor[rgb]{0.9,0.9,0.9}{
            \makecell[{{p{\linewidth}}}]{
                \texttt{\tiny{[GM$\langle$P2]}}
                \texttt{MOVERR: L, (2,2)} \\
\texttt{MOVERR: P, (2,3)} \\
            }
        }
    }
     \\ \\

    \theutterance \stepcounter{utterance}  
    & & & \multicolumn{2}{p{0.3\linewidth}}{
        \cellcolor[rgb]{0.9,0.9,0.9}{
            \makecell[{{p{\linewidth}}}]{
                \texttt{\tiny{[GM$|$GM]}}
                \texttt{MOVERR: L, (2,2)} \\
\texttt{MOVERR: P, (2,3)} \\
            }
        }
    }
    & & \\ \\

    \theutterance \stepcounter{utterance}  
    & & & \multicolumn{2}{p{0.3\linewidth}}{
        \cellcolor[rgb]{0.9,0.9,0.9}{
            \makecell[{{p{\linewidth}}}]{
                \texttt{\tiny{[GM$|$GM]}}
                \texttt{Invalid response format} \\
            }
        }
    }
    & & \\ \\

    \theutterance \stepcounter{utterance}  
    & & \multicolumn{4}{p{0.6\linewidth}}{
        \cellcolor[rgb]{0.9,0.9,0.9}{
            \makecell[{{p{\linewidth}}}]{
                \texttt{\tiny{[GM$\rangle$P2]}}
                \texttt{Penalización: ¡Tu mensaje no tiene el formato esperado!} \\
\\ 
\texttt{Asegúrate de que tu respuesta solo contenga `DECIR: <MENSAJE>` o `MOVERR: <OBJETO>, (<X>, <Y>)`, ¡y nada más!} \\
\\ 
\texttt{Han acumulado colectivamente 1 de 8 penalizaciones.} \\
\texttt{¡Por favor, inténtalo de nuevo!} \\
            }
        }
    }
    & \\ \\

    \theutterance \stepcounter{utterance}  
    & & & \multicolumn{4}{p{0.6\linewidth}}{
        \cellcolor[rgb]{0.9,0.9,0.9}{
            \makecell[{{p{\linewidth}}}]{
                \texttt{\tiny{[GM$\langle$P2]}}
                \texttt{DECIR: Confirmo que realizaré los movimientos para mis 'L' y 'P' hacia (2,2) y (2,3) respectivamente, ya que ambas celdas están vacías según tu confirmación.} \\
            }
        }
    }
     \\ \\

    \theutterance \stepcounter{utterance}  
    & & & \multicolumn{2}{p{0.3\linewidth}}{
        \cellcolor[rgb]{0.9,0.9,0.9}{
            \makecell[{{p{\linewidth}}}]{
                \texttt{\tiny{[GM$|$GM]}}
                \texttt{DECIR: Confirmo que realizaré los movimientos para mis 'L' y 'P' hacia (2,2) y (2,3) respectivamente, ya que ambas celdas están vacías según tu confirmación.} \\
            }
        }
    }
    & & \\ \\

    \theutterance \stepcounter{utterance}  
    & & \multicolumn{4}{p{0.6\linewidth}}{
        \cellcolor[rgb]{0.9,0.9,0.9}{
            \makecell[{{p{\linewidth}}}]{
                \texttt{\tiny{[P1$\langle$GM]}}
                \texttt{Tu mensaje ha sido transmitido al otro jugador.} \\
\\ 
\\ 
\texttt{Actualmente estás jugando la ronda 3 de un máximo de 12.} \\
\texttt{Han acumulado colectivamente 1 de 8 penalizaciones.} \\
\texttt{El otro jugador envió este mensaje:} \\
\\ 
\\ 
\texttt{"Confirmo que realizaré los movimientos para mis 'L' y 'P' hacia (2,2) y (2,3) respectivamente, ya que ambas celdas están vacías según tu confirmación."} \\
\\ 
\\ 
\texttt{¿Cuál es tu próximo comando?} \\
            }
        }
    }
    & \\ \\

    \theutterance \stepcounter{utterance}  
    & \multicolumn{4}{p{0.6\linewidth}}{
        \cellcolor[rgb]{0.9,0.9,0.9}{
            \makecell[{{p{\linewidth}}}]{
                \texttt{\tiny{[P1$\rangle$GM]}}
                \texttt{MOVERR: L, (2,2)} \\
            }
        }
    }
    & & \\ \\

    \theutterance \stepcounter{utterance}  
    & & & \multicolumn{2}{p{0.3\linewidth}}{
        \cellcolor[rgb]{0.9,0.9,0.9}{
            \makecell[{{p{\linewidth}}}]{
                \texttt{\tiny{[GM$|$GM]}}
                \texttt{MOVERR: L, (2,2)} \\
            }
        }
    }
    & & \\ \\

    \theutterance \stepcounter{utterance}  
    & & & \multicolumn{2}{p{0.3\linewidth}}{
        \cellcolor[rgb]{0.9,0.9,0.9}{
            \makecell[{{p{\linewidth}}}]{
                \texttt{\tiny{[GM$|$GM]}}
                \texttt{Invalid response format} \\
            }
        }
    }
    & & \\ \\

    \theutterance \stepcounter{utterance}  
    & & \multicolumn{4}{p{0.6\linewidth}}{
        \cellcolor[rgb]{0.9,0.9,0.9}{
            \makecell[{{p{\linewidth}}}]{
                \texttt{\tiny{[P1$\langle$GM]}}
                \texttt{Penalización: ¡Tu mensaje no tiene el formato esperado!} \\
\\ 
\texttt{Asegúrate de que tu respuesta solo contenga `DECIR: <MENSAJE>` o `MOVERR: <OBJETO>, (<X>, <Y>)`, ¡y nada más!} \\
\\ 
\texttt{Han acumulado colectivamente 2 de 8 penalizaciones.} \\
\texttt{¡Por favor, inténtalo de nuevo!} \\
            }
        }
    }
    & \\ \\

    \theutterance \stepcounter{utterance}  
    & \multicolumn{4}{p{0.6\linewidth}}{
        \cellcolor[rgb]{0.9,0.9,0.9}{
            \makecell[{{p{\linewidth}}}]{
                \texttt{\tiny{[P1$\rangle$GM]}}
                \texttt{DECIR: Confirmo que mi 'L' está en (2,2) y la celda está vacía, así que puedo moverla allí. ¿Podrías ejecutar el movimiento MOVERR: L, (2,2) en tu tablero?} \\
            }
        }
    }
    & & \\ \\

    \theutterance \stepcounter{utterance}  
    & & & \multicolumn{2}{p{0.3\linewidth}}{
        \cellcolor[rgb]{0.9,0.9,0.9}{
            \makecell[{{p{\linewidth}}}]{
                \texttt{\tiny{[GM$|$GM]}}
                \texttt{DECIR: Confirmo que mi 'L' está en (2,2) y la celda está vacía, así que puedo moverla allí. ¿Podrías ejecutar el movimiento MOVERR: L, (2,2) en tu tablero?} \\
            }
        }
    }
    & & \\ \\

    \theutterance \stepcounter{utterance}  
    & & \multicolumn{4}{p{0.6\linewidth}}{
        \cellcolor[rgb]{0.9,0.9,0.9}{
            \makecell[{{p{\linewidth}}}]{
                \texttt{\tiny{[GM$\rangle$P2]}}
                \texttt{Tu mensaje ha sido transmitido al otro jugador.} \\
\\ 
\\ 
\texttt{Actualmente estás jugando la ronda 4 de un máximo de 12.} \\
\texttt{Han acumulado colectivamente 2 de 8 penalizaciones.} \\
\texttt{El otro jugador envió este mensaje:} \\
\\ 
\\ 
\texttt{"Confirmo que mi 'L' está en (2,2) y la celda está vacía, así que puedo moverla allí. ¿Podrías ejecutar el movimiento MOVERR: L, (2,2) en tu tablero?"} \\
\\ 
\\ 
\texttt{¿Cuál es tu próximo comando?} \\
            }
        }
    }
    & \\ \\

    \theutterance \stepcounter{utterance}  
    & & & \multicolumn{4}{p{0.6\linewidth}}{
        \cellcolor[rgb]{0.9,0.9,0.9}{
            \makecell[{{p{\linewidth}}}]{
                \texttt{\tiny{[GM$\langle$P2]}}
                \texttt{MOVERR: L, (2, 2)} \\
            }
        }
    }
     \\ \\

    \theutterance \stepcounter{utterance}  
    & & & \multicolumn{2}{p{0.3\linewidth}}{
        \cellcolor[rgb]{0.9,0.9,0.9}{
            \makecell[{{p{\linewidth}}}]{
                \texttt{\tiny{[GM$|$GM]}}
                \texttt{MOVERR: L, (2, 2)} \\
            }
        }
    }
    & & \\ \\

    \theutterance \stepcounter{utterance}  
    & & & \multicolumn{2}{p{0.3\linewidth}}{
        \cellcolor[rgb]{0.9,0.9,0.9}{
            \makecell[{{p{\linewidth}}}]{
                \texttt{\tiny{[GM$|$GM]}}
                \texttt{Invalid response format} \\
            }
        }
    }
    & & \\ \\

    \theutterance \stepcounter{utterance}  
    & & \multicolumn{4}{p{0.6\linewidth}}{
        \cellcolor[rgb]{0.9,0.9,0.9}{
            \makecell[{{p{\linewidth}}}]{
                \texttt{\tiny{[GM$\rangle$P2]}}
                \texttt{Penalización: ¡Tu mensaje no tiene el formato esperado!} \\
\\ 
\texttt{Asegúrate de que tu respuesta solo contenga `DECIR: <MENSAJE>` o `MOVERR: <OBJETO>, (<X>, <Y>)`, ¡y nada más!} \\
\\ 
\texttt{Han acumulado colectivamente 3 de 8 penalizaciones.} \\
\texttt{¡Por favor, inténtalo de nuevo!} \\
            }
        }
    }
    & \\ \\

    \theutterance \stepcounter{utterance}  
    & & & \multicolumn{4}{p{0.6\linewidth}}{
        \cellcolor[rgb]{0.9,0.9,0.9}{
            \makecell[{{p{\linewidth}}}]{
                \texttt{\tiny{[GM$\langle$P2]}}
                \texttt{MOVERR: L, (2,2)} \\
            }
        }
    }
     \\ \\

    \theutterance \stepcounter{utterance}  
    & & & \multicolumn{2}{p{0.3\linewidth}}{
        \cellcolor[rgb]{0.9,0.9,0.9}{
            \makecell[{{p{\linewidth}}}]{
                \texttt{\tiny{[GM$|$GM]}}
                \texttt{MOVERR: L, (2,2)} \\
            }
        }
    }
    & & \\ \\

    \theutterance \stepcounter{utterance}  
    & & & \multicolumn{2}{p{0.3\linewidth}}{
        \cellcolor[rgb]{0.9,0.9,0.9}{
            \makecell[{{p{\linewidth}}}]{
                \texttt{\tiny{[GM$|$GM]}}
                \texttt{Invalid response format} \\
            }
        }
    }
    & & \\ \\

    \theutterance \stepcounter{utterance}  
    & & \multicolumn{4}{p{0.6\linewidth}}{
        \cellcolor[rgb]{0.9,0.9,0.9}{
            \makecell[{{p{\linewidth}}}]{
                \texttt{\tiny{[GM$\rangle$P2]}}
                \texttt{Penalización: ¡Tu mensaje no tiene el formato esperado!} \\
\\ 
\texttt{Asegúrate de que tu respuesta solo contenga `DECIR: <MENSAJE>` o `MOVERR: <OBJETO>, (<X>, <Y>)`, ¡y nada más!} \\
\\ 
\texttt{Han acumulado colectivamente 4 de 8 penalizaciones.} \\
\texttt{¡Por favor, inténtalo de nuevo!} \\
            }
        }
    }
    & \\ \\

    \theutterance \stepcounter{utterance}  
    & & & \multicolumn{4}{p{0.6\linewidth}}{
        \cellcolor[rgb]{0.9,0.9,0.9}{
            \makecell[{{p{\linewidth}}}]{
                \texttt{\tiny{[GM$\langle$P2]}}
                \texttt{MOVERR: L, (2, 2)} \\
            }
        }
    }
     \\ \\

    \theutterance \stepcounter{utterance}  
    & & & \multicolumn{2}{p{0.3\linewidth}}{
        \cellcolor[rgb]{0.9,0.9,0.9}{
            \makecell[{{p{\linewidth}}}]{
                \texttt{\tiny{[GM$|$GM]}}
                \texttt{MOVERR: L, (2, 2)} \\
            }
        }
    }
    & & \\ \\

    \theutterance \stepcounter{utterance}  
    & & & \multicolumn{2}{p{0.3\linewidth}}{
        \cellcolor[rgb]{0.9,0.9,0.9}{
            \makecell[{{p{\linewidth}}}]{
                \texttt{\tiny{[GM$|$GM]}}
                \texttt{Invalid response format} \\
            }
        }
    }
    & & \\ \\

    \theutterance \stepcounter{utterance}  
    & & \multicolumn{4}{p{0.6\linewidth}}{
        \cellcolor[rgb]{0.9,0.9,0.9}{
            \makecell[{{p{\linewidth}}}]{
                \texttt{\tiny{[GM$\rangle$P2]}}
                \texttt{Penalización: ¡Tu mensaje no tiene el formato esperado!} \\
\\ 
\texttt{Asegúrate de que tu respuesta solo contenga `DECIR: <MENSAJE>` o `MOVERR: <OBJETO>, (<X>, <Y>)`, ¡y nada más!} \\
\\ 
\texttt{Han acumulado colectivamente 5 de 8 penalizaciones.} \\
\texttt{¡Por favor, inténtalo de nuevo!} \\
            }
        }
    }
    & \\ \\

    \theutterance \stepcounter{utterance}  
    & & & \multicolumn{4}{p{0.6\linewidth}}{
        \cellcolor[rgb]{0.9,0.9,0.9}{
            \makecell[{{p{\linewidth}}}]{
                \texttt{\tiny{[GM$\langle$P2]}}
                \texttt{MOVERR: L, (2, 2)} \\
            }
        }
    }
     \\ \\

    \theutterance \stepcounter{utterance}  
    & & & \multicolumn{2}{p{0.3\linewidth}}{
        \cellcolor[rgb]{0.9,0.9,0.9}{
            \makecell[{{p{\linewidth}}}]{
                \texttt{\tiny{[GM$|$GM]}}
                \texttt{MOVERR: L, (2, 2)} \\
            }
        }
    }
    & & \\ \\

    \theutterance \stepcounter{utterance}  
    & & & \multicolumn{2}{p{0.3\linewidth}}{
        \cellcolor[rgb]{0.9,0.9,0.9}{
            \makecell[{{p{\linewidth}}}]{
                \texttt{\tiny{[GM$|$GM]}}
                \texttt{Invalid response format} \\
            }
        }
    }
    & & \\ \\

    \theutterance \stepcounter{utterance}  
    & & \multicolumn{4}{p{0.6\linewidth}}{
        \cellcolor[rgb]{0.9,0.9,0.9}{
            \makecell[{{p{\linewidth}}}]{
                \texttt{\tiny{[GM$\rangle$P2]}}
                \texttt{Penalización: ¡Tu mensaje no tiene el formato esperado!} \\
\\ 
\texttt{Asegúrate de que tu respuesta solo contenga `DECIR: <MENSAJE>` o `MOVERR: <OBJETO>, (<X>, <Y>)`, ¡y nada más!} \\
\\ 
\texttt{Han acumulado colectivamente 6 de 8 penalizaciones.} \\
\texttt{¡Por favor, inténtalo de nuevo!} \\
            }
        }
    }
    & \\ \\

    \theutterance \stepcounter{utterance}  
    & & & \multicolumn{4}{p{0.6\linewidth}}{
        \cellcolor[rgb]{0.9,0.9,0.9}{
            \makecell[{{p{\linewidth}}}]{
                \texttt{\tiny{[GM$\langle$P2]}}
                \texttt{MOVERR: L, (2, 2)} \\
            }
        }
    }
     \\ \\

    \theutterance \stepcounter{utterance}  
    & & & \multicolumn{2}{p{0.3\linewidth}}{
        \cellcolor[rgb]{0.9,0.9,0.9}{
            \makecell[{{p{\linewidth}}}]{
                \texttt{\tiny{[GM$|$GM]}}
                \texttt{MOVERR: L, (2, 2)} \\
            }
        }
    }
    & & \\ \\

    \theutterance \stepcounter{utterance}  
    & & & \multicolumn{2}{p{0.3\linewidth}}{
        \cellcolor[rgb]{0.9,0.9,0.9}{
            \makecell[{{p{\linewidth}}}]{
                \texttt{\tiny{[GM$|$GM]}}
                \texttt{Invalid response format} \\
            }
        }
    }
    & & \\ \\

    \theutterance \stepcounter{utterance}  
    & & \multicolumn{4}{p{0.6\linewidth}}{
        \cellcolor[rgb]{0.9,0.9,0.9}{
            \makecell[{{p{\linewidth}}}]{
                \texttt{\tiny{[GM$\rangle$P2]}}
                \texttt{Penalización: ¡Tu mensaje no tiene el formato esperado!} \\
\\ 
\texttt{Asegúrate de que tu respuesta solo contenga `DECIR: <MENSAJE>` o `MOVERR: <OBJETO>, (<X>, <Y>)`, ¡y nada más!} \\
\\ 
\texttt{Han acumulado colectivamente 7 de 8 penalizaciones.} \\
\texttt{¡Por favor, inténtalo de nuevo!} \\
            }
        }
    }
    & \\ \\

    \theutterance \stepcounter{utterance}  
    & & & \multicolumn{4}{p{0.6\linewidth}}{
        \cellcolor[rgb]{0.9,0.9,0.9}{
            \makecell[{{p{\linewidth}}}]{
                \texttt{\tiny{[GM$\langle$P2]}}
                \texttt{MOVERR: L, (2, 2)} \\
            }
        }
    }
     \\ \\

    \theutterance \stepcounter{utterance}  
    & & & \multicolumn{2}{p{0.3\linewidth}}{
        \cellcolor[rgb]{0.9,0.9,0.9}{
            \makecell[{{p{\linewidth}}}]{
                \texttt{\tiny{[GM$|$GM]}}
                \texttt{MOVERR: L, (2, 2)} \\
            }
        }
    }
    & & \\ \\

    \theutterance \stepcounter{utterance}  
    & & & \multicolumn{2}{p{0.3\linewidth}}{
        \cellcolor[rgb]{0.9,0.9,0.9}{
            \makecell[{{p{\linewidth}}}]{
                \texttt{\tiny{[GM$|$GM]}}
                \texttt{Invalid response format} \\
            }
        }
    }
    & & \\ \\

    \theutterance \stepcounter{utterance}  
    & & \multicolumn{4}{p{0.6\linewidth}}{
        \cellcolor[rgb]{0.9,0.9,0.9}{
            \makecell[{{p{\linewidth}}}]{
                \texttt{\tiny{[GM$\rangle$P2]}}
                \texttt{Penalización: ¡Tu mensaje no tiene el formato esperado!} \\
\\ 
\texttt{Asegúrate de que tu respuesta solo contenga `DECIR: <MENSAJE>` o `MOVERR: <OBJETO>, (<X>, <Y>)`, ¡y nada más!} \\
\\ 
\texttt{Han acumulado colectivamente 8 de 8 penalizaciones.} \\
\texttt{¡Por favor, inténtalo de nuevo!} \\
            }
        }
    }
    & \\ \\

    \theutterance \stepcounter{utterance}  
    & & & \multicolumn{4}{p{0.6\linewidth}}{
        \cellcolor[rgb]{0.9,0.9,0.9}{
            \makecell[{{p{\linewidth}}}]{
                \texttt{\tiny{[GM$\langle$P2]}}
                \texttt{No puedo continuar con el juego bajo estas condiciones de formato. Adjunto mi mensaje en el formato requerido:} \\
\\ 
\texttt{DECIR: He recibido múltiples penalizaciones por no seguir el formato correcto. Intento corregir mi siguiente comando.} \\
\\ 
\texttt{MOVERR: L, (2, 2)} \\
            }
        }
    }
     \\ \\

    \theutterance \stepcounter{utterance}  
    & & & \multicolumn{2}{p{0.3\linewidth}}{
        \cellcolor[rgb]{0.9,0.9,0.9}{
            \makecell[{{p{\linewidth}}}]{
                \texttt{\tiny{[GM$|$GM]}}
                \texttt{No puedo continuar con el juego bajo estas condiciones de formato. Adjunto mi mensaje en el formato requerido:} \\
\\ 
\texttt{DECIR: He recibido múltiples penalizaciones por no seguir el formato correcto. Intento corregir mi siguiente comando.} \\
\\ 
\texttt{MOVERR: L, (2, 2)} \\
            }
        }
    }
    & & \\ \\

    \theutterance \stepcounter{utterance}  
    & & & \multicolumn{2}{p{0.3\linewidth}}{
        \cellcolor[rgb]{0.9,0.9,0.9}{
            \makecell[{{p{\linewidth}}}]{
                \texttt{\tiny{[GM$|$GM]}}
                \texttt{Invalid format: head and tail are not empty} \\
\texttt{head: 'No puedo continuar con el juego bajo estas condiciones de formato. Adjunto mi mensaje en el formato requerido:} \\
\\ 
\texttt{'} \\
\texttt{tail: '} \\
\\ 
\texttt{MOVERR: L, (2, 2)'} \\
            }
        }
    }
    & & \\ \\

    \theutterance \stepcounter{utterance}  
    & & & \multicolumn{2}{p{0.3\linewidth}}{
        \cellcolor[rgb]{0.9,0.9,0.9}{
            \makecell[{{p{\linewidth}}}]{
                \texttt{\tiny{[GM$|$GM]}}
                \texttt{Maximum number of penalties exceeded} \\
            }
        }
    }
    & & \\ \\

    \theutterance \stepcounter{utterance}  
    & & & \multicolumn{2}{p{0.3\linewidth}}{
        \cellcolor[rgb]{0.9,0.9,0.9}{
            \makecell[{{p{\linewidth}}}]{
                \texttt{\tiny{[GM$|$GM]}}
                \texttt{Initial states:} \\
\texttt{```} \\
\texttt{Player 1:} \\
\texttt{ 1234567} \\
\texttt{╔═╤═════╗} \\
\texttt{║◌│◌◌◌◌L║ 1} \\
\texttt{║◌│◌◌◌◌◌║ 2} \\
\texttt{║◌│◌◌◌◌◌║ 3} \\
\texttt{║◌│◌◌◌◌◌║ 4} \\
\texttt{╟─┼───┐C║ 5} \\
\texttt{║◌│◌◌◌├─╢ 6} \\
\texttt{║P│◌◌◌│◌║ 7} \\
\texttt{╚═╧═══╧═╝} \\
\\ 
\\ 
\texttt{Player 2:} \\
\texttt{ 1234567} \\
\texttt{╔═╤═════╗} \\
\texttt{║◌│◌◌◌◌◌║ 1} \\
\texttt{║◌│◌◌PL◌║ 2} \\
\texttt{║◌│◌◌◌◌◌║ 3} \\
\texttt{║◌│◌◌◌◌◌║ 4} \\
\texttt{╟─┼───┐◌║ 5} \\
\texttt{║◌│◌◌◌├─╢ 6} \\
\texttt{║◌│◌C◌│◌║ 7} \\
\texttt{╚═╧═══╧═╝} \\
\\ 
\texttt{```} \\
            }
        }
    }
    & & \\ \\

    \theutterance \stepcounter{utterance}  
    & & & \multicolumn{2}{p{0.3\linewidth}}{
        \cellcolor[rgb]{0.9,0.9,0.9}{
            \makecell[{{p{\linewidth}}}]{
                \texttt{\tiny{[GM$|$GM]}}
                \texttt{End states:} \\
\texttt{```} \\
\texttt{Player 1:} \\
\texttt{ 1234567} \\
\texttt{╔═╤═════╗} \\
\texttt{║◌│◌◌◌◌L║ 1} \\
\texttt{║◌│◌◌◌◌◌║ 2} \\
\texttt{║◌│◌◌◌◌◌║ 3} \\
\texttt{║◌│◌◌◌◌◌║ 4} \\
\texttt{╟─┼───┐C║ 5} \\
\texttt{║◌│◌◌◌├─╢ 6} \\
\texttt{║P│◌◌◌│◌║ 7} \\
\texttt{╚═╧═══╧═╝} \\
\\ 
\\ 
\texttt{Player 2:} \\
\texttt{ 1234567} \\
\texttt{╔═╤═════╗} \\
\texttt{║◌│◌◌◌◌◌║ 1} \\
\texttt{║◌│◌◌PL◌║ 2} \\
\texttt{║◌│◌◌◌◌◌║ 3} \\
\texttt{║◌│◌◌◌◌◌║ 4} \\
\texttt{╟─┼───┐◌║ 5} \\
\texttt{║◌│◌◌◌├─╢ 6} \\
\texttt{║◌│◌C◌│◌║ 7} \\
\texttt{╚═╧═══╧═╝} \\
\\ 
\texttt{```} \\
            }
        }
    }
    & & \\ \\

    \theutterance \stepcounter{utterance}  
    & & & \multicolumn{2}{p{0.3\linewidth}}{
        \cellcolor[rgb]{0.9,0.9,0.9}{
            \makecell[{{p{\linewidth}}}]{
                \texttt{\tiny{[GM$|$GM]}}
                \texttt{* success: False} \\
\texttt{* lose: False} \\
\texttt{* aborted: True} \\
\texttt{{-}{-}{-}{-}{-}{-}{-}} \\
\texttt{* Shifts: 0.00} \\
\texttt{* Max Shifts: 4.00} \\
\texttt{* Min Shifts: 2.00} \\
\texttt{* End Distance Sum: 11.42} \\
\texttt{* Init Distance Sum: 11.42} \\
\texttt{* Expected Distance Sum: 12.57} \\
\texttt{* Penalties: 9.00} \\
\texttt{* Max Penalties: 8.00} \\
\texttt{* Rounds: 3.00} \\
\texttt{* Max Rounds: 12.00} \\
\texttt{* Object Count: 3.00} \\
            }
        }
    }
    & & \\ \\

\end{supertabular}
}

\end{document}
